\section{Security Analysis and Discussion}
\label{sec:analysis}

\subsection{Scope and Goal}

This section analyzes the parts of Ghostext that are already implemented and directly measurable. It is not a complete indistinguishability proof against adaptive steganalysis. The goal is narrower: isolate each approximation source in the current pipeline, express its effect with explicit quantities, and define what remains to be measured beyond the released pilot baseline.

\subsection{Distribution Pipeline Decomposition}

At token step $t$, the implementation transforms the base next-token distribution $P_t$ through three deterministic operators:
\begin{equation}
P_t \xrightarrow{\text{truncate}} \widetilde{P}_t \xrightarrow{\text{quantize}} Q_t \xrightarrow{\text{stability filter}} Q'_t.
\end{equation}

The truncation operator keeps a top-$p$/max-$k$ subset and removes tail mass. Quantization maps real probabilities to integer frequencies with total budget $F_{\mathrm{tot}}$. Stability filtering removes candidates that fail render/re-tokenize consistency and then renormalizes.

This decomposition is useful because each term is observable in logs and can be ablated independently.

Configuration note: the released real-backend baseline uses $F_{\mathrm{tot}}=4096$ and $h_{\min}=0.0$, while implementation defaults are $F_{\mathrm{tot}}=65536$ and $h_{\min}=1.0$. All numeric examples below are labeled explicitly, and the released efficiency numbers should be interpreted as an ungated initial operating point rather than a recommended deployment setting.

\subsection{Per-Step Total-Variation Bound}

Section~\ref{sec:protocol-details} defines
\begin{equation}
\alpha_t = 1-\sum_{v\in C_t} P_t(v),
\end{equation}
\begin{equation}
\beta_t = \sum_{v\in C_t\setminus S_t} Q_t(v),
\end{equation}
and provides
\begin{equation}
\mathrm{TV}(Q'_t,P_t)\le\alpha_t + \frac{N_t}{2F_{\mathrm{tot}}} + \beta_t,
\end{equation}
where $N_t=|C_t|$.

The three addends have clear meanings.
$\alpha_t$ is policy truncation mass,
$N_t/(2F_{\mathrm{tot}})$ is finite-frequency quantization distortion,
and $\beta_t$ is additional mass removed by tokenization-stability filtering, including the normalization effect after filtering.

\subsection{Quantization Bias and a Conservative KL Upper Bound}

To answer the reviewer concern on quantization bias, define
$\delta_t(v)=Q_t(v)-\widetilde{P}_t(v)$.
By construction, $\sum_v\delta_t(v)=0$ and $|\delta_t(v)|\le 1/F_{\mathrm{tot}}$.
Because every quantized candidate gets at least frequency one, $Q_t(v)\ge 1/F_{\mathrm{tot}}$.

Using $\widetilde{P}_t(v)=Q_t(v)-\delta_t(v)$:
\begin{align}
D_{\mathrm{KL}}(\widetilde{P}_t\|Q_t)
&=\sum_{v:\widetilde{P}_t(v)>0} \widetilde{P}_t(v)\log\!\left(\frac{\widetilde{P}_t(v)}{Q_t(v)}\right).
\end{align}
Applying the standard $D_{\mathrm{KL}}(P\|Q)\le \chi^2(P,Q)$ upper bound~\cite{tsybakov2009introduction} gives
\begin{align}
D_{\mathrm{KL}}(\widetilde{P}_t\|Q_t)
&\le \chi^2(\widetilde{P}_t,Q_t)
= \sum_{v:\widetilde{P}_t(v)>0} \frac{(\widetilde{P}_t(v)-Q_t(v))^2}{Q_t(v)}
= \sum_{v:\widetilde{P}_t(v)>0} \frac{\delta_t(v)^2}{Q_t(v)}.
\end{align}
Then, using $Q_t(v)\ge 1/F_{\mathrm{tot}}$ and $|\delta_t(v)|\le 1/F_{\mathrm{tot}}$,
\begin{align}
D_{\mathrm{KL}}(\widetilde{P}_t\|Q_t)
&\le \sum_v \frac{(1/F_{\mathrm{tot}})^2}{1/F_{\mathrm{tot}}}
\le \frac{N_t}{F_{\mathrm{tot}}}\quad\text{(nats)}.
\end{align}

The resulting bound is loose but explicit. Under implementation defaults ($N_t\le16$, $F_{\mathrm{tot}}=65536$), it is at most $16/65536 \approx 2.44\times10^{-4}$ nats for quantization alone. Under the released baseline setting ($N_t\le64$, $F_{\mathrm{tot}}=4096$), the corresponding worst-case ceiling is $64/4096 = 1.56\times10^{-2}$ nats. Truncation and stability filtering can dominate this term in practice, so all terms should be logged together.

\subsection{Entropy-Threshold Gating and Capacity}

Embedding is allowed only when both conditions hold: $N_t\ge2$ and per-step entropy is at least $h_{\min}$.
The implementation default is $h_{\min}=1.0$ bit, while the released real-backend baseline sets $h_{\min}=0.0$ to avoid entropy-threshold gating during this initial recoverability/runtime run.

Let $I_t\in\{0,1\}$ indicate whether step $t$ allows embedding, and let $\Delta b_t$ denote interval-resolved bits at that step.
An implementation-level effective payload rate can be summarized as
\begin{equation}
\eta_{\mathrm{pkt}}=\frac{\sum_t I_t\,\Delta b_t}{\text{packet-carrying tokens}}.
\end{equation}

The threshold therefore trades capacity for synchronization safety.
Lower $h_{\min}$ increases opportunities to embed, but also increases exposure to low-entropy stalls and retry events. Accordingly, the released $2.478$ bits/token figure should be read as an ungated initial baseline under $h_{\min}=0.0$, not as a safe-default deployment recommendation.
The practical next step is therefore not to headline this single number, but to run the same condition grid at least over $h_{\min}\in\{0.0,0.5,1.0\}$ before treating any efficiency figure as representative of a conservative deployment envelope.

\subsection{Sequence-Level Accumulation and Diagnostics}

For full-sequence distributions, chain decomposition gives
\begin{equation}
D_{\mathrm{KL}}(Q_{1:T}\|P_{1:T})
=\sum_{t=1}^{T}\mathbb{E}_{Q_{<t}}\left[D_{\mathrm{KL}}(Q_t(\cdot\mid Y_{<t})\|P_t(\cdot\mid Y_{<t}))\right].
\end{equation}

In Ghostext, contexts depend on prior stego choices, so a tight analytic bound is difficult without stronger assumptions. The practical path is to log per-step terms and compute empirical diagnostics per message:
\begin{equation}
\widehat{B}_{\mathrm{TV}}=\sum_{t=1}^{T}\left(\alpha_t + \frac{N_t}{2F_{\mathrm{tot}}} + \beta_t\right).
\label{eq:bhat-tv}
\end{equation}

If $P_{1:T}$ denotes the matched-distribution natural-cover process, then sequence-level distinguishability is governed by $\mathrm{TV}(Q_{1:T},P_{1:T})$. With $\mathrm{Adv}_{\mathrm{censor}} := |2\Pr[\mathrm{correct}] - 1|$ for a binary hypothesis test between one sample from $Q_{1:T}$ and one sample from $P_{1:T}$, an idealized censor satisfies
\begin{equation}
\mathrm{Adv}_{\mathrm{censor}}\le \mathrm{TV}(Q_{1:T},P_{1:T}) \le \sqrt{\tfrac{1}{2}D_{\mathrm{KL}}(Q_{1:T}\|P_{1:T})},
\end{equation}
where the first inequality is the standard testing bound and the second is Pinsker's inequality~\cite{tsybakov2009introduction}. The steganographic meaning of that bridge is then interpreted through frameworks such as Cachin~\cite{cachin1998information} and Hopper et al.~\cite{hopper2002provably}. That sequence-level bridge is the quantity a censor-facing claim would ultimately need.

$\widehat{B}_{\mathrm{TV}}$ is therefore an audit signal rather than a direct detector-risk guarantee. In this revision we use it conservatively in two ways: a large value is direct evidence that approximation terms are not negligible, while a small value is only suggestive and still requires detector-side validation. Its practical usefulness is therefore calibration-dependent: if a later regime shows small $\widehat{B}_{\mathrm{TV}}$ but high detector AUC, the metric should be downgraded to a local implementation-health diagnostic rather than treated as a censor-risk proxy. For that reason, the staged detector block treats correlation analysis between $\widehat{B}_{\mathrm{TV}}$ and detector scores as a required output, not an optional add-on. The full real-backend measurement block remains pending (Stage 3).

\subsection{Relation to Provable-Security Limits}

A useful limit case is $\alpha_t\to0$ (no truncation), $\beta_t\to0$ (no stability filtering loss), and $F_{\mathrm{tot}}\to\infty$ (no quantization loss), under exact sender-receiver replay of the same next-token distribution. In this limit, the implementation approaches ideal arithmetic-coding steganography over the model distribution at each step.

This limit clarifies where Ghostext sits relative to Cachin-style distribution matching and Hopper-style stego security formulations: our current system is an approximation to that ideal regime, with approximation terms made explicit by $\alpha_t$, $N_t/(2F_{\mathrm{tot}})$, and $\beta_t$. The present paper does not claim that current finite-precision, truncated, stability-filtered operation already satisfies a full provable-security theorem.

\subsection{Fingerprint Collision and Misbinding}

The synchronization fingerprint is a truncated 64-bit SHA-256 digest over runtime configuration, backend metadata, and prompt hash. For $n$ honest sessions,
\begin{equation}
\Pr[\text{collision}]\approx\frac{n(n-1)}{2^{65}}.
\end{equation}
For $n=10^6$, this is roughly $2.7\times10^{-8}$.

A collision is not sufficient for successful forgery. Decoding still requires authenticated packet validity under the shared passphrase and protocol checks. The fingerprint should be interpreted as a fast mismatch detector, not as a stand-alone authenticity primitive.
Under an idealized uniform-target approximation, $2^{40}$ adaptive collision attempts against a fixed 64-bit target still yield only about $2^{-24}$ success probability, which is small but not a margin worth presenting as future-proof when the byte overhead of widening is negligible. The fingerprint is also a Bob-local synchronization check rather than a public plaintext header field in the passive model of this paper. The 64-bit width is therefore a deployment tradeoff in the current artifact, not a claim that 64 bits is universally optimal. Given the negligible byte overhead of widening this field, 128-bit fingerprinting is the conservative future direction unless compatibility with already released artifacts must be preserved.

\subsection{Retry-Induced Traffic Leakage Model}

Retries improve robustness in difficult regions but can change externally observable traffic.
Let $R$ be attempt count, $L$ output length, and $T_{\mathrm{lat}}$ latency.
Observed marginals become mixtures:
\begin{equation}
\Pr[L=\ell]=\sum_r \Pr[R=r]\Pr[L=\ell\mid R=r],
\end{equation}
\begin{equation}
\Pr[T_{\mathrm{lat}}\le t]=\sum_r \Pr[R=r]\Pr[T_{\mathrm{lat}}\le t\mid R=r].
\end{equation}

Importantly, $R$ is sender-local: failed attempts are discarded and Bob only decodes the final transmitted cover. So $R$ does not alter decode replay correctness, but it can affect sender-side timing/energy and external traffic observables.

Even when text-only quality is strong, these mixture shifts may be detectable by a traffic observer. In the released 18-trial pilot, \texttt{attempts\_used} has mean $1.778$ and maximum $4$: 11 runs finished in one attempt, 2 in two attempts, 3 in three attempts, and 2 in four attempts. The English subset is mostly single-attempt (7 of 9), while the Chinese subset spreads across one, two, three, and four attempts with counts 4, 1, 2, and 2. By contrast, a natural-cover generator without a retry loop has degenerate $R=1$ unless dummy retries are added deliberately. Under the current policy, retry count alone is therefore a trivially separating feature against any zero-retry natural-cover baseline; no learned detector is required to exploit that difference. This already makes the side-channel concern concrete in the released pilot, and it also explains why the released p99 encode latency is driven by a few high-retry cases rather than a stabilized tail law. Evaluation should therefore report retry-conditioned and unconditional length/latency distributions, and we keep this as a pending detector/traffic block beyond the current baseline.

\subsection{Prompt-Knowledge and Multi-Message Regimes}

The released real-backend evidence is a template-known operational floor: prompts are fixed public templates repeated across cases. For broader deployment claims, two stronger regimes still need direct measurement:
(1) private-prompt or unknown-prompt conditions, where the censor does not know the sender's exact prompt,
and (2) sample-accumulation, where many messages from one sender are available.

If a detector score is $s(x)$ per message and the censor aggregates
$\bar{s}_m = \frac{1}{m}\sum_{i=1}^{m}s(x_i)$,
small per-message effects can become detectable as $m$ grows.
This is why single-message and multi-message tests should both be first-class reporting targets.

\subsection{Synchronization Assumptions in Deployment}

Current evidence is strongest for local, pinned deployments where both parties control model files, tokenizer behavior, backend build, and runtime policy. The implementation binds backend metadata and tokenizer hash into the configuration fingerprint, which is effective for detecting drift in this deployment mode.

When using third-party hosted APIs, exact backend-version observability may be limited. In that environment, strict deterministic replay can be unavailable by construction, and Ghostext should be treated as out-of-scope unless an API provider exposes stable and verifiable model/tokenizer/version identifiers sufficient for replay.

\subsection{Cross-Hardware and Numeric Determinism Caveat}

The current fingerprint includes backend metadata exposed by the backend implementation, but it does not include a full hardware fingerprint (for example, GPU model and driver stack). In local llama.cpp CPU mode this is less exposed to GPU-kernel nondeterminism than many hosted GPU pipelines, but floating-point and platform differences can still perturb logits near candidate cut boundaries.

The practical implication is conservative: cross-hardware deployment should be validated empirically before operational use, and hardware/runtime context should be logged as a reproducibility field in evaluation artifacts. The current paper therefore claims deterministic replay only within a pinned local environment, not as a universal cross-hardware property.

\subsection{Natural Tail Policy Clarification}

Current natural-tail tokens are sampled from the same constrained candidate mechanism used by packet-carrying tokens, then excluded from decode once packet recovery is complete. This avoids a second sampling pathway but may preserve policy artifacts in tail regions.

The correct interpretation is therefore operational, not security-complete:
in this version, natural tails mainly extend text to a more complete ending, but they are not a distinct detector-evasion mechanism. They should be measured separately from packet-carrying prefixes in both text and traffic analyses.

\subsection{Logging Schema for Measurable Analysis}

\begin{table}[t]
\caption{Per-step and per-message diagnostics and metadata to log for analysis, ablation, and reproducibility.}
\label{tab:analysis-logs}
\begin{tabular}{p{0.25\linewidth}p{0.68\linewidth}}
\toprule
Field & Role in analysis \\
\midrule
$\alpha_t$ & Truncation mass removed from base model distribution at step $t$. \\
$N_t$ & Candidate count after truncation and before quantization. \\
$F_{\mathrm{tot}}$ & Integer frequency budget used by quantizer. \\
$\beta_t$ & Quantized mass removed by retokenization-stability filtering. \\
$H_t$ & Candidate entropy used for min-entropy gating and low-entropy window checks. \\
$I_t$ & Embedding-enabled indicator at step $t$ (1 if payload bits may be consumed). \\
$R$ & Retry count per message; needed for latency/length mixture analysis. \\
Failure class & One of explicit fail-closed classes (integrity, mismatch, stall, budget, etc.). \\
Model metadata & Model family/version and tokenizer hash used in generation. \\
Policy config & top-$p$, max-$k$, min-entropy threshold, total frequency budget, retry limits. \\
Prompt regime & Prompt family label and whether template-known or private-prompt setting. \\
Runtime context & Hardware class, backend build info, and software versions. \\
\bottomrule
\end{tabular}
\end{table}

\subsection{Current Claims Only}

At this stage, supported claims are restricted to deterministic decode under matched runtime state, explicit mismatch detection, fail-closed behavior, and measurable decomposition of approximation sources.

Claims about broad detector resistance and cross-family generalization remain pending dedicated detector and cross-family experiments. Active text modification in transit remains out of scope in this revision, and deterministic claims should be read as scoped to the pinned local environment used for the released artifact.
